\documentclass[11pt]{article}

\usepackage[margin=1in]{geometry}
\usepackage{amsmath,amssymb,amsthm}
\usepackage{graphicx}
\usepackage{enumitem}
\usepackage[hidelinks]{hyperref}

\newtheorem{theorem}{Theorem}
\newtheorem{lemma}[theorem]{Lemma}
\newcommand{\dist}{\operatorname{dist}}
\newcommand{\vol}{\operatorname{vol}}

\title{The Growing-Degree Threshold for\\
Geometric Embeddings of Trees}
\author{FA/Banach run \texttt{fa\_banach\_001}}
\date{9 August 2026}

\begin{document}
\maketitle

\section*{Status}

This packet records a candidate full answer to Problem~1.6 of
Altschuler--Dodos--Tikhomirov--Tyros \cite{ADTT}. The proof combines their
random embedding and local-lemma architecture with the
Brascamp--Lieb--Luttinger rearrangement inequality \cite{BLL}. The current
verdict is
\[
 \boxed{\textbf{likely valid full result; human review requested.}}
\]
The result determines the threshold up to universal constant factors for the
entire range $3\leq\Delta\leq N-1$. It does not optimize the constants and
does not address the separate optimal-constants Problem~1.5 of \cite{ADTT}.

\section*{Source problem}

A graph $G=(V,E)$ is \emph{geometrically embeddable} into a normed space $X$
if there is a map $\zeta:V\to X$ such that, for all distinct $u,v\in V$,
\[
 \|\zeta(u)-\zeta(v)\|_X\leq1
 \quad\Longleftrightarrow\quad \{u,v\}\in E.
\]
The source proves a universal $\Theta(\bigl(\log N\bigr)/\log\log N)$
threshold when the maximum degree $\Delta$ is fixed. On PDF page~3, in
Problem~1.6, it asks for the threshold when $\Delta$ grows with $N$.

\begin{center}
\includegraphics[width=\textwidth]{figures/open_problem_crop.png}
\end{center}

In transcription, Problem~1.6 asks for an explicit function $m(N,\Delta)$
and universal constants $\alpha_0,\alpha_1$ such that every $N$-vertex tree
of maximum degree $\Delta$ embeds into every normed space of dimension at
least $\alpha_1m$, while a complete tree of degree $\Delta$ does not embed
below $\alpha_0m$.

\section*{Proof Intuition}

Put
\[
 H=\frac{\log N}{\log(\Delta-1)}.
\]
This is the height scale of a complete rooted degree-$\Delta$ tree with $N$
vertices. Packing the images of a bipartite independent set in such a tree
forces dimension of order
\[
 q(N,\Delta)=\frac{\log N}{\log(2+H)}.
\]

For the matching upper bound, the source attaches independent uniform vectors
from the unit ball of the target norm to the tree edges, then adds a small
independent Gaussian regularization. Its proof already works at dimension
$Aq$ for every path except the finitely many shortest lengths. The published
thin-shell estimate controls such a short path only by
$\exp(-m^{\kappa/2})$, which cannot pay for a dependency neighborhood of size
$\Delta^{O(1)}$ when $\Delta$ is large.

The missing estimate is supplied by rearrangement. Conditional on all but two
edge vectors, a short-path failure asks whether the sum of two independent
uniform points of an arbitrary symmetric convex body lies in a translate of a
slightly enlarged copy of that body. The Brascamp--Lieb--Luttinger inequality
says that this probability is maximized by a Euclidean ball. In the Euclidean
ball, radii are exponentially concentrated near the boundary, while two
independent directions have exponentially small probability of making the
fixed obtuse angle required for their sum to have norm near one. Consequently
the short-path probability is $O(e^{-cm})$, uniformly over every norm.

This exponential estimate pays for $\Delta^{O(1)}$ in the sparse-dependency
regime. If $\Delta$ is so large that the local lemma uses its dense $N^2$
dependency bound instead, then $H=O(1)$ and $q\asymp\log N$, so the same
estimate pays for $N^2$. All remaining path lengths are controlled by the
source's $\exp(-cm\log k)$ estimate and its Gaussian completion.

\section*{The full theorem}

\begin{theorem}[Growing-degree threshold]\label{thm:main}
There are universal constants $A<\infty$ and $N_0\in\mathbb N$ with the
following property. Let $N\geq N_0$ and $3\leq\Delta\leq N-1$, and define
\begin{equation}\label{eq:qdef}
 q(N,\Delta)=
 \frac{\log N}{
 \log\!\left(2+\dfrac{\log N}{\log(\Delta-1)}\right)}.
\end{equation}
Then:
\begin{enumerate}[label=(\roman*)]
 \item every tree on $N$ vertices with maximum degree at most $\Delta$
 geometrically embeds into every normed space of dimension at least
 $\lceil A q(N,\Delta)\rceil$;
 \item the breadth-first truncation to $N$ vertices of the complete rooted
 degree-$\Delta$ tree does not geometrically embed into any normed space of
 dimension less than $q(N,\Delta)/4$.
\end{enumerate}
At the cardinalities of full complete rooted trees, the witness in (ii) is
the full complete tree. Thus Problem~1.6 is answered by the explicit threshold
function $m(N,\Delta)=q(N,\Delta)$, with universal constants
$\alpha_0=1/4$ and $\alpha_1=A$.
\end{theorem}

\section*{Proof}

We first prove the norm-uniform estimate that removes the only fixed-degree
step from the source proof.

\begin{lemma}[Exponential short-path estimate]\label{lem:short}
Let $\kappa\in(0,1)$ be the absolute exponent in Theorem~2.2 of \cite{ADTT}.
There are absolute constants $c,C>0$ and $m_\ast\in\mathbb N$ such that the
following holds. Let $X=(\mathbb R^m,\|\cdot\|_X)$ be any normed space with
$m\geq m_\ast$, and let $Y_1,\ldots,Y_k$ be independent random vectors
uniformly distributed in its unit ball, where $k\geq2$. Then
\begin{equation}\label{eq:short}
 \mathbb P\!\left(
  \left\|\sum_{i=1}^kY_i\right\|_X
  \leq1+\frac{1}{2m^{1-\kappa}}
 \right)\leq Ce^{-cm}.
\end{equation}
The same estimate holds with arbitrary fixed signs in the sum.
\end{lemma}

\begin{proof}
Write $K=B_X$ and
\[
 t=1+\frac{1}{2m^{1-\kappa}}.
\]
The body $K$ is centrally symmetric. After increasing $m_\ast$, we may
assume $t\leq9/8$.

Condition on $Z=Y_3+\cdots+Y_k$ (with $Z=0$ if $k=2$). For every realization
$Z=z$, the conditional probability in \eqref{eq:short} is
\begin{equation}\label{eq:threeint}
 \frac{1}{\vol(K)^2}
 \int_{\mathbb R^m}\!\int_{\mathbb R^m}
 \mathbf 1_K(x)\mathbf 1_K(y)
 \mathbf 1_{tK-z}(x+y)\,dx\,dy.
\end{equation}
Let $B$ be the centered Euclidean ball having the same volume as $K$. The
Riesz rearrangement inequality, equivalently the three-function case of the
Brascamp--Lieb--Luttinger inequality \cite{BLL}, bounds
\eqref{eq:threeint} by
\[
 \frac{1}{\vol(B)^2}
 \int_{\mathbb R^m}\!\int_{\mathbb R^m}
 \mathbf 1_B(x)\mathbf 1_B(y)\mathbf 1_{tB}(x+y)\,dx\,dy.
\]
Here the symmetric decreasing rearrangements of $K,K,tK-z$ are respectively
$B,B,tB$. By scaling, the last display is
\begin{equation}\label{eq:euclprob}
 \mathbb P\bigl(\|U+V\|_2\leq t\bigr),
\end{equation}
where $U,V$ are independent and uniform in the Euclidean unit ball.

Write $U=R\Theta$ and $V=S\Phi$, where the directions are independent and
uniform on $S^{m-1}$, independent also of the radii. Since
\[
 \mathbb P(R<15/16)=\mathbb P(S<15/16)=(15/16)^m,
\]
the probability that either radius is smaller than $15/16$ is at most
$2(15/16)^m$. On the complementary event, if
$\|R\Theta+S\Phi\|_2\leq9/8$, then
\[
 \langle\Theta,\Phi\rangle
 \leq
 \frac{(9/8)^2-R^2-S^2}{2RS}
 \leq
 \frac{(9/8)^2-2(15/16)^2}{2(15/16)^2}
 =-\frac7{25}.
\]
The middle expression is decreasing in each of $R,S\in[15/16,1]$.

For a fixed direction, the first coordinate of a uniform point on
$S^{m-1}$ has density
\[
 c_m(1-s^2)^{(m-3)/2}\mathbf1_{[-1,1]}(s),
 \qquad
 c_m=\frac{\Gamma(m/2)}{\sqrt\pi\,\Gamma((m-1)/2)}\leq\sqrt m.
\]
It follows, for an absolute $c_0>0$, that
\[
 \mathbb P\bigl(\langle\Theta,\Phi\rangle\leq-7/25\bigr)
 \leq \sqrt m\left(1-\frac{49}{625}\right)^{(m-3)/2}
 \leq C_0e^{-c_0m}.
\]
Together with the radial estimate, this bounds \eqref{eq:euclprob} by
$Ce^{-cm}$, uniformly in $z$ and $k$. Averaging over $Z$ proves
\eqref{eq:short}. Fixed signs do not alter the distribution because the
$Y_i$ are symmetric.
\end{proof}

We now prove Theorem~\ref{thm:main}. Put
\[
 D=\log(\Delta-1),\qquad
 H=\frac{\log N}{D},\qquad
 L=\log(2+H),\qquad
 q=\frac{\log N}{L}.
\]
Because $3\leq\Delta\leq N-1$, we have $H\geq1$. Uniformly over this entire
range,
\begin{equation}\label{eq:qrange}
 c\frac{\log N}{\log\log N}\leq q\leq C\log N.
\end{equation}
Also, since $H\mapsto\log(2+H)/H$ is decreasing on $[1,\infty)$,
\begin{equation}\label{eq:Dq}
 D=\frac{\log N}{H}\leq(\log3)q.
\end{equation}

Retain the source parameters
\[
 c_1=\frac14,
 \qquad
 c_2=\min\left\{\frac14,\frac\kappa4\right\},
 \qquad
 \ell_0=\exp\bigl((\log N)^{c_2}\bigr),
\]
and its absolute integer $k_0\geq3$ from (4.7). Set
\[
 m=\lceil Aq\rceil,
 \qquad
 k_1=\lfloor H\rfloor,
\]
where the absolute constant $A$ will be chosen sufficiently large. From
\eqref{eq:qrange}, uniformly in $\Delta$,
\begin{equation}\label{eq:small}
 \log\ell_0=(\log N)^{c_2}=o(q),
 \qquad
 k_1<\ell_0
\end{equation}
for all sufficiently large $N$.

We use the independent uniform edge vectors, path events, dependency graph,
and random embedding from Sections~4.1--4.3 of \cite{ADTT}. For a path of
length $k$, let the uniform-vector bad event be
\[
 \mathcal L_{u,v}=
 \begin{cases}
 \left\{\left\|\sum_{e\in P(u,v)}\!\!\pm Y_e\right\|_X
 \leq1+\dfrac{1}{2m^{1-\kappa}}\right\},&2\leq k\leq k_0,\\[8pt]
 \left\{\left\|\sum_{e\in P(u,v)}\!\!\pm Y_e\right\|_X
 \leq k^{c_1}\right\},&k_0<k\leq\ell_0.
 \end{cases}
\]
The signs depend on the position of the root, but do not change the
distribution. Assign every length-$k$ event the local-lemma weight
\begin{equation}\label{eq:bk}
 b_k=
 \begin{cases}
 \displaystyle\frac{1}{2^{k+1}6\ell_0k(\Delta-1)^k},
 &2\leq k\leq k_1,\\[7pt]
 \displaystyle\frac{3}{2\pi^2k^2N^2},
 &k_1<k\leq\ell_0.
 \end{cases}
\end{equation}
The first range is allowed to be empty.

For $k\leq k_1$, $(\Delta-1)^k\leq N$. Since
$k\leq H=O(\log N)$ and $\ell_0=N^{o(1)}$, uniformly
\begin{equation}\label{eq:dependency}
 6k\ell_0(\Delta-1)^k\leq N^2
\end{equation}
for all sufficiently large $N$. Claim~4.1 and the calculation in (4.16) of
\cite{ADTT} therefore give, for every path-event $\mathcal L_{u,v}$,
\begin{equation}\label{eq:product}
 \prod_{\mathcal L_{u',v'}\sim\mathcal L_{u,v}}
 (1-b_{\dist(u',v')})\geq\frac12.
\end{equation}

It remains to show
\begin{equation}\label{eq:compare}
 \mathbb P(\mathcal L_{u,v})\leq\frac{b_k}{2}
 \qquad (2\leq k\leq\ell_0).
\end{equation}
We record three elementary comparisons. Their constants depend only on the
absolute integer $k_0$:
\begin{align}
 k_0<k\leq k_1
 &\quad\Longrightarrow\quad
 k\log(2\Delta-2)\leq Cq\log k,
 \label{eq:sparsecost}\\
 k>k_0\ \hbox{and}\ k>k_1
 &\quad\Longrightarrow\quad
 q\log k\geq c\log N,
 \label{eq:densegain}\\
 2\leq k\leq k_0\ \hbox{and}\ k>k_1
 &\quad\Longrightarrow\quad
 q\geq c\log N.
 \label{eq:shortdense}
\end{align}
Indeed, for \eqref{eq:sparsecost}, monotonicity of $x/\log x$ gives
\[
 \frac{kD}{\log k}\leq\frac{HD}{\log H}
 =\frac{\log N}{\log H}\leq Cq,
\]
because $H\geq k>k_0$ and $L/\log H$ is then bounded; also
$\log(2\Delta-2)\leq2D$. For \eqref{eq:densegain}, if $H\geq k_0$, then
$k>k_1$ implies $k>H$ and $\log k/L\geq c$; if $H<k_0$, use
$k>k_0$ and $L\leq\log(k_0+2)$. Finally, \eqref{eq:shortdense} follows
because $k>k_1$ implies $H<k\leq k_0$, so $L$ is bounded.

Suppose first that $2\leq k\leq k_0$. Lemma~\ref{lem:short} gives
\begin{equation}\label{eq:shortprob}
 \mathbb P(\mathcal L_{u,v})\leq Ce^{-cm}\leq Ce^{-cAq}.
\end{equation}
If $k\leq k_1$, then \eqref{eq:Dq}, \eqref{eq:small}, and
\eqref{eq:bk} give
\[
 \log\frac2{b_k}\leq Cq.
\]
If $k>k_1$, then \eqref{eq:shortdense} gives the same conclusion from
$\log(2/b_k)\leq2\log N+o(\log N)$. Thus \eqref{eq:compare} follows in
the short range once $A$ is a sufficiently large absolute constant.

Next assume $k_0<k\leq k_1$. Lemma~3.2 of \cite{ADTT}, with $c_1=1/4$,
gives
\begin{equation}\label{eq:midprob}
 \mathbb P(\mathcal L_{u,v})
 \leq\exp\left(-\frac m8\log k\right).
\end{equation}
By \eqref{eq:sparsecost}, \eqref{eq:small}, and \eqref{eq:bk},
\[
 \log\frac2{b_k}\leq Cq\log k.
\]
Choosing $A>8C$ proves \eqref{eq:compare} in this range.

Finally assume $k>k_0$ and $k>k_1$. The estimate \eqref{eq:midprob} still
holds, while \eqref{eq:densegain} yields
\[
 \frac m8\log k\geq cA\log N.
\]
On the other hand,
\[
 \log\frac2{b_k}
 =2\log N+2\log k+O(1)
 =(2+o(1))\log N,
\]
because $k\leq\ell_0=N^{o(1)}$. A final universal enlargement of $A$
proves \eqref{eq:compare}. The cases are exhaustive.

Combining \eqref{eq:compare} with \eqref{eq:product}, the asymmetric
Lovasz local lemma gives positive probability to the event $\mathcal L$ on
which every nonadjacent pair at tree distance at most $\ell_0$ has the
required uniform-vector separation.

It remains to check that the Gaussian completion in Section~4.3 of
\cite{ADTT} is uniform in $\Delta$. From \eqref{eq:qrange}, our dimension
satisfies
\begin{equation}\label{eq:mrange}
 c_A\frac{\log N}{\log\log N}\leq m\leq C_A\log N.
\end{equation}
The source uses
\[
 \delta=
 \exp\left(-\left(\frac12-\frac{c_1}{2}\right)(\log N)^{c_2}\right).
\]
For adjacent pairs, for distances $2\leq k\leq k_0$, and for distances
$k_0<k\leq\ell_0$, its Gaussian upper-tail thresholds contain respectively
the factors
\[
 \frac{\delta^{-1}}{\sqrt m\log N},
 \qquad
 \frac{\delta^{-1}}{m^{3/2-\kappa}},
 \qquad
 \frac{\exp((c_1/2)(\log N)^{c_2})}{\sqrt m}.
\]
The exponential in $(\log N)^{c_2}$ dominates every displayed power of
$m$ and $\log N$ uniformly on \eqref{eq:mrange}. Equations (4.29)--(4.31)
of \cite{ADTT} therefore bound each corresponding conditional bad
probability by $N^{-2}$. For $k>\ell_0$, equations (4.32)--(4.35) give the
bound
\[
 \exp\bigl(-c\,m(\log N)^{c_2}\bigr)\leq N^{-2},
\]
again uniformly by \eqref{eq:mrange}. These estimates use only the Gaussian
variables and are unchanged after conditioning on $\mathcal L$.

There are fewer than $N^2/2$ unordered vertex pairs. Conditioned on
$\mathcal L$, the union bound is strictly less than $1/2$, so the random map
is a geometric embedding with positive probability. If the ambient normed
space has dimension larger than $m$, carry out the construction in any
$m$-dimensional subspace. This proves part~(i).

For part~(ii), use the breadth-first truncated complete tree from part~(A) of
Theorem~1.1 of \cite{ADTT}. It has $N$ vertices, maximum degree at most
$\Delta$, a root $r$, and radius
\[
 h_0=\left\lceil
 \frac{\log\left((N-1)(\Delta-2)/\Delta+1\right)}
 {\log(\Delta-1)}
 \right\rceil\leq H+1.
\]
Suppose it geometrically embeds into a $d$-dimensional normed space. Since
the tree is bipartite, it has an independent set $V_0$ with
$|V_0|\geq N/2$. The radius-$1/2$ balls centered at the images of vertices
of $V_0$ are pairwise disjoint, and all lie in the radius-$(h_0+1/2)$ ball
about the image of $r$. Comparing Lebesgue volumes gives
\[
 \frac N2\leq(2h_0+1)^d\leq(2H+3)^d.
\]
For $N\geq4$, using $2H+3\leq(H+2)^2$, we obtain
\[
 d\geq\frac{\log(N/2)}{\log(2H+3)}
 \geq\frac{\log N}{4\log(H+2)}=\frac q4.
\]
Thus no such embedding exists in dimension less than $q/4$. At full
complete-tree cardinalities no truncation occurs. This proves part~(ii) and
the theorem. \hfill$\Box$

\section*{Verification notes}

The proof was audited against Sections~3 and~4 of the source. The new input is
only Lemma~\ref{lem:short}; it replaces Lemma~3.3 of the source in the finite
range $2\leq k\leq k_0$. The event threshold is identical, so no later part
of the random embedding changes.

The degree-dependent checks are:
\begin{enumerate}
 \item the sparse path-dependency inequality \eqref{eq:dependency};
 \item the short sparse comparison, paid for by \eqref{eq:Dq};
 \item the short dense comparison, where $H=O(1)$ forces
 $q\asymp\log N$;
 \item the transition at $k_1=\lfloor H\rfloor$, covered by
 \eqref{eq:sparsecost}--\eqref{eq:shortdense}.
\end{enumerate}
All Gaussian estimates are degree-free and use only the uniform dimension
window \eqref{eq:mrange}. No numerical experiment is used as proof.

The main human-review focus should be the application of rearrangement to the
translated third body in \eqref{eq:threeint}, and the four exhaustive
local-lemma subranges (short/sparse, short/dense, long/sparse, long/dense).

\section*{Novelty search and limitations}

A bounded literature search on 9 August 2026 checked arXiv:2504.15212, its
May 2026 Combinatorica publication, the exact title, the exact phrase
``Problem 1.6 (Trees of growing maximal degree)'', the authors, and keyword
combinations involving geometric embeddings, trees, growing maximum degree,
normed spaces, rearrangement, and kissing probability. It found the source
paper and the classical rearrangement theorem, but no separate paper claiming
to answer Problem~1.6 or giving the threshold \eqref{eq:qdef}. Novelty
confidence is moderate rather than definitive because the source is recent.

The proof uses a classical theorem in a new place; it does not claim the
rearrangement inequality itself as new. The constant $A$ is universal but is
not numerically optimized, since the source constant $k_0$ already depends on
implicit absolute constants in the slicing estimates. The result also does
not settle the optimal universal constants asked in Problem~1.5.

\section*{Human-review recommendation}

Review Lemma~\ref{lem:short} first. If its translated-set rearrangement step
and spherical-cap estimate pass, check the four local-lemma comparisons and
the assertion that Section~4.3 of the source is uniform on
\eqref{eq:mrange}. If those checks pass, this packet gives a complete
constant-factor answer to Problem~1.6 for all growing degrees.

\small
\begin{thebibliography}{9}
\bibitem{ADTT}
D.~J. Altschuler, P.~Dodos, K.~Tikhomirov, and K.~Tyros,
\emph{A universal threshold for geometric embeddings of trees},
Combinatorica 46 (2026), article 20; arXiv:2504.15212v2,
\url{https://arxiv.org/abs/2504.15212}.

\bibitem{BLL}
H.~J. Brascamp, E.~H. Lieb, and J.~M. Luttinger,
\emph{A general rearrangement inequality for multiple integrals},
J. Funct. Anal. 17 (1974), 227--237,
\url{https://doi.org/10.1016/0022-1236(74)90013-5}.
\end{thebibliography}

\end{document}
